\section{Results}

\begin{frame}{The ANN baseline gives a stable IV surrogate}
  \begin{columns}[T,onlytextwidth]
    \begin{column}{0.67\textwidth}
      \centering
      \includegraphics[width=\linewidth,height=0.78\textheight,keepaspectratio]{nn/sensitivity_grid_3x2_log.png}
    \end{column}
    \begin{column}{0.29\textwidth}
      \begin{itemize}
        \item Test IV RMSE:
        \[
          1.86\cdot 10^{-3}.
        \]
        \item Stable behaviour in the interpolation core.
        \item Error increases near boundary and extrapolation regions.
      \end{itemize}
    \end{column}
  \end{columns}
\end{frame}

\begin{frame}{Sobolev training concentrates ANN Greek errors}
  \centering
  \includegraphics[width=\linewidth,height=0.73\textheight,keepaspectratio]{nn/ann_iv_baseline_vs_sobolev_error_hist_overlay_logy_5greeks.png}

  \vspace{1mm}
  \begin{itemize}
    \item Sobolev training narrows the error distributions for Delta, Gamma, Theta, and Rho.
    \item Vega also improves, although with a smaller separation between variants.
  \end{itemize}
\end{frame}

\begin{frame}{Sobolev improves ANN Greeks across the grid}
  \centering
  \includegraphics[width=\linewidth,height=0.76\textheight,keepaspectratio]{ann_sobolev_greek_rel_error_maps_landscape.png}

  \vspace{1mm}
  \begin{center}
    \scriptsize
    Greek RMSE ratios: Delta \(7.0\times\), Gamma \(12.7\times\), Vega \(2.2\times\), Theta \(5.0\times\), Rho \(20.8\times\).
  \end{center}
\end{frame}

\begin{frame}{The Sobolev IV surrogate gives tighter CaNN residuals}
  \begin{columns}[T,onlytextwidth]
    \begin{column}{0.50\textwidth}
      \centering
      \includegraphics[width=\linewidth,height=0.67\textheight,keepaspectratio]{cann/cann_residual_heatmap_m_tau_l2_reg_smooth.png}
    \end{column}
    \begin{column}{0.46\textwidth}
      \centering
      \resizebox{\linewidth}{!}{%
      \begin{tabular}{lccc}
        \toprule
        Forward model & MSE & RMSE & \(\max |e|\) \\
        \midrule
        ANN baseline & \(7.43\cdot10^{-10}\) & \(2.73\cdot10^{-5}\) & \(6.79\cdot10^{-5}\) \\
        ANN Sobolev & \(3.69\cdot10^{-11}\) & \(6.08\cdot10^{-6}\) & \(1.36\cdot10^{-5}\) \\
        Liu et al. & \(\sim10^{-8}\) & -- & -- \\
        \bottomrule
      \end{tabular}}

      \vspace{6mm}
      \begin{itemize}
        \item Calibration is performed on the replicated \(35\)-quote Liu grid.
        \item The Sobolev forward surrogate gives the tightest IV residuals.
        \item Parameter recovery improves for four of the five Heston parameters relative to the baseline CaNN.
      \end{itemize}
    \end{column}
  \end{columns}
\end{frame}

\begin{frame}{The baseline PINN is strongest as a pricing model}
  \begin{columns}[T,onlytextwidth]
    \begin{column}{0.70\textwidth}
      \centering
      \includegraphics[width=\linewidth,height=0.78\textheight,keepaspectratio]{pinn/pricing/pricing_benchmark_composite.png}
    \end{column}
    \begin{column}{0.26\textwidth}
      \begin{itemize}
        \item Fixed grid:
        \[
          n=25{,}921.
        \]
        \item Price RMSE:
        \[
          3.07\cdot10^{-4}.
        \]
        \item MAPE:
        \[
          0.28\%.
        \]
      \end{itemize}
    \end{column}
  \end{columns}
\end{frame}

\begin{frame}{Sobolev improves PINN Greeks under a different regime}
  \centering
  \small
  \begin{tabular}{lccc}
    \toprule
    Greek & PINN baseline RMSE & PINN Sobolev RMSE & Ratio \\
    \midrule
    Delta & \(3.33\cdot10^{-2}\) & \(3.51\cdot10^{-4}\) & \(94.9\times\) \\
    Gamma & \(3.85\cdot10^{-1}\) & \(8.94\cdot10^{-3}\) & \(43.1\times\) \\
    Vega  & \(6.55\cdot10^{-2}\) & \(3.17\cdot10^{-3}\) & \(20.7\times\) \\
    Theta & \(4.82\cdot10^{-3}\) & \(5.30\cdot10^{-4}\) & \(9.1\times\) \\
    Rho   & \(8.37\cdot10^{-2}\) & \(1.43\cdot10^{-3}\) & \(58.5\times\) \\
    \bottomrule
  \end{tabular}

  \vspace{6mm}
  \begin{itemize}
    \item Sobolev terms strongly improve full-grid Greek diagnostics.
    \item The interpretation is not a uniform model improvement: derivative labels change the information regime.
    \item Greek throughput changes from \(9.48\cdot10^3\) to \(8.60\cdot10^3\) points/s.
  \end{itemize}
\end{frame}

\begin{frame}{Short maturities remain the hardest PINN region}
  \centering
  \includegraphics[width=\linewidth,height=0.76\textheight,keepaspectratio]{pinn_sobolev_greek_rel_error_maps_landscape.png}

  \vspace{1mm}
  \begin{itemize}
    \item Sobolev reduces global errors, but the short-maturity region remains structured.
    \item This motivates a separate local diagnostic around the payoff kink.
  \end{itemize}
\end{frame}

\begin{frame}{ACV targets the short-maturity ATM hard region}
  \centering
  \includegraphics[width=\linewidth,height=0.77\textheight,keepaspectratio]{pinn_sobolev_acv_greek_rel_error_maps_landscape.png}

  \vspace{1mm}
  \begin{center}
    \scriptsize
    Hard-region grid: \(m\in[0.9704,1.0305]\), \(\tau\in[0.005,0.05]\).
  \end{center}
\end{frame}

\begin{frame}{Neural evaluators are faster after training}
  \centering
  \begin{tabular}{lrr}
    \toprule
    Component & Throughput & Unit \\
    \midrule
    ANN evaluator & \(8.43\cdot10^{5}\) & IV evaluations/s \\
    PINN pricing & \(4.91\cdot10^{5}\) & price points/s \\
    COS repeated pricing & \(8.736\cdot10^{1}\) & price points/s \\
    PINN Greek baseline & \(9.48\cdot10^{3}\) & Greek points/s \\
    PINN Greek Sobolev & \(8.60\cdot10^{3}\) & Greek points/s \\
    \bottomrule
  \end{tabular}

  \vspace{7mm}
  \begin{itemize}
    \item The speed advantage appears after training, in repeated evaluation mode.
    \item The comparison is not used to rank all models globally; it quantifies the computational role of each branch.
  \end{itemize}
\end{frame}
